\documentclass[10pt,a4paper]{article}
\usepackage{tikz} % for drawing figures
\usepackage{amsmath} % for equations
\usepackage{url} % for URLs
\usepackage{graphicx}
\usepackage{multicol}
\usepackage{varwidth}
\usepackage{blindtext}


\usepackage{linguex} % ** special include in directory: for doing handy example labeling and bracketing
\renewcommand{\firstrefdash}{} % used for linguex package not to put hyphens in example refs (1a instead of 1-a)
\usepackage{cogsci}
\usepackage{pslatex}
\usepackage{apacite}

\newcommand{\sem}[1]{\mbox{$[\![$#1$]\!]$}}
\newcommand{\lam}{$\lambda$}
\newcommand{\gcs}[1]{\textcolor{blue}{[gcs: #1]}} 



\title{On the added informativity of ambiguous language}
\author{\large \textbf{our names}\\
our emails\\
our affiliations}


\begin{document}
\maketitle

\begin{abstract}
XXX


\textbf{Keywords:} 
ambiguity, pragmatics, informativity, Rational Speech Act models

\end{abstract}

\section{Introduction}

Linguists often treat ambiguity as a bug in the communication system, something to be avoided or explained away \cite{grice1975,chomsky2002minimalism}. However, recent research has begun to take notice of the efficiency ambiguity affords us: by relying on context to fill in missing information, we can reuse lightweight bits of language rather than fully specifying the intended message \cite{levinson2000,piantadosietal2012,wasow2015}. In other words, ambiguity serves as a feature---not a bug---of an efficient communication system. This reasoning accords with years of psycholinguistic research documenting that speakers readily produce ambiguous utterances (see \citeNP{ferreira2008}, for an overview). 

The current work identifies an additional benefit afforded by using ambiguous language: the \emph{extra} information we gain from observing how our listeners resolve ambiguity. We advance the hypothesis that language users learn about each other's private knowledge by observing how they resolve ambiguity. If language does not do the job of specifying the information necessary for full interpretation, then listeners are left to draw on their beliefs and preferences to fill in the gaps; by observing how listeners fill those gaps in, speakers learn about the beliefs and preferences the listeners are using. 

By way of illustration, take, for example, the scenario in Fig.~\ref{FG-ref-game}. Suppose a speaker produces the single-word utterance ``blue'' in an attempt to signal one of the objects to a listener. The utterance is ambiguous; it can pick out either the blue square or the blue circle. Suppose further that, upon hearing ``blue,'' the listener selects the blue circle. In observing this choice, the speaker stands to learn something about the private thoughts of the listener: what made her select the blue circle instead of the blue square? Perhaps the circle is more salient to the listener, or the listener has a preference for circles, or the listener believes the speaker has a preference for circles. By observing how the listener resolves the ambiguity in reference, the listener can learn this information. 

However, accessing this added information requires the speaker to reason pragmatically about the pragmatic reasoning of the listener---a higher-order pragmatic reasoning, as it were.

\begin{figure}
	\centering
	\includegraphics[width=2in]{images/rsa_scene.eps}
	\caption{A simple reference game scenario from \citeA{frankgoodman2012}. In the game, speakers choose a single-word utterance to signal one of the objects to a listener. In this scenario, the speaker chooses between the utterances ``blue,'' ``green,'' ``square,'' and ``circle.''}\label{FG-ref-game}
	\end{figure}


\section{Background}

\cite{wasow2015}

\cite{piantadosietal2012}

\cite{frankgoodman2012,goodmanfrank2016}


\section{Model}


\section{Expt.~1: Inferring preferences}


\section{Expt.~2: Choosing utterances}


\section{Discussion}



\bibliographystyle{apacite}
\setlength{\bibleftmargin}{.125in}
\setlength{\bibindent}{-\bibleftmargin}

\bibliography{prior-inference}

\end{document}

